\vfill%
\pagebreak
\section{Pattern matching}
\label{sec:flow:pattern_matching}

\token{match} compares a value with a sequence of patterns, and evaluates the
body of the first one that fits. A pattern tests the value, and can also take it
apart, declaring variables for its parts. Patterns are also used by
\token{if let} (Section~\ref{sec:flow:cond_pattern}), \token{while let}
(Section~\ref{sec:flow:while_let_loop}), variable declarations
(Section~\ref{sec:memory:pattern_vdecl}) and \token{catch}
(Chapter~\ref{chap:error}); this section defines them for all five.

\subsection{Grammar}

\begin{lstlisting}[style=grammarVerb]
match_expression := 'match' expression '{' arm+ '}'
arm              := pattern ('if' condition)? '=>' body
body             := block
                  | expression ';'?

pattern := '_'
         | binding
         | tuple_pattern
         | array_pattern
         | option_pattern
         | object_pattern
         | expression ('|' expression)*

binding        := 'ref'? ('mut' | 'dmut')? Identifier
                  (':' (type | '_'))? ('=' pattern)?
tuple_pattern  := '(' ')'
                | '(' pattern ',' ')'
                | '(' pattern (',' pattern)+ ')'
                | '(' pattern (',' pattern)* '...' ')'
array_pattern  := '[' ']'
                | '[' pattern (',' pattern)+ ']'
                | '[' pattern (',' pattern)* '...' ']'
option_pattern := ('Ok' | 'Err') '(' pattern? ')'
object_pattern := type '(' named_patterns? ')'
named_patterns := named_pattern (',' named_pattern)*
named_pattern  := Identifier '->' pattern
\end{lstlisting}

The arms follow one another with no separator. The guard of an arm is introduced
by \token{if}, where \token{if let} and \token{while let} use \token{\&\&}.

\subsubsection{Bodies that are single expressions}

A body that is not a block is read as the longest expression possible. When the
next arm starts with a token that can continue an expression -- an opening
parenthesis or bracket, or a binary operator such as \token{-} -- the body
continues into it. In the next listing, the arm \token{(0, \_) => 0} followed by
the pattern \token{(x, y)} reads as the call \token{0(x, y)}, and the match is a
syntax error; a body \token{"zero"} followed by an arm \token{-1 => ...} would
read as \token{"zero" - 1}. A body followed by such an arm must be a block, or
end with a semicolon.

\begin{lstlisting}[style=coloredverbatimError, escapechar=@]
let s = match t {
  (0, _) => @\hb{0}@ // read as 0(x, y) => ...
  (x, y) => x + y
};
\end{lstlisting}

\subsubsection{Brackets around a single pattern}

Parentheses around a single pattern with no comma group it, and do not form a
tuple: \token{(x)} is the binding \token{x}, and \token{(x,)} the tuple of one
field. Square brackets behave the same way: \token{[x]} is the binding
\token{x}, not an array of one element, and there is no pattern for an array of
exactly one element.

\subsection{Evaluation}

The matched value is evaluated once. The arms are then tried in order. An arm is
selected when its pattern fits the value and its guard, if any, holds; the guard
is evaluated only when the pattern fits, with the variables of the pattern in
scope. The body of the selected arm is evaluated, and no further arm is tried. If
no arm is selected, nothing is evaluated.

\subsection{Patterns}

\begin{center}
  \begin{adjustbox}{max width=\linewidth}
    \begin{tabular}{|l|l|l|}
      \hline
      Pattern & Fits the value \token{v} when & Declares\\[0pt]
      \hline
      \hline
      \token{\_} & always & nothing\\[0pt]
      \token{e}, an expression & \token{v == e} & nothing\\[0pt]
      \token{a .. b}, \token{a ... b}, with \token{v} an integer & \token{v in a .. b}, \token{v in a ... b} & nothing\\[0pt]
      \token{e1 | e2 | ...} & one of the alternatives fits & nothing\\[0pt]
      \token{x} & always & \token{x}\\[0pt]
      \token{x: T} & \token{v} is of type \token{T} & \token{x}, of type \token{T}\\[0pt]
      \token{x: T = p} & \token{v} is of type \token{T} and fits \token{p} & \token{x}, and the variables of \token{p}\\[0pt]
      \token{(p1, ..., pn)} & \token{v} is a tuple of $n$ fields, each fitting its pattern & the variables of each \token{pi}\\[0pt]
      \token{[p1, ..., pn]} & \token{v} has $n$ elements, each fitting its pattern & the variables of each \token{pi}\\[0pt]
      \token{Ok(p)} & \token{v} is an option holding a value that fits \token{p} & the variables of \token{p}\\[0pt]
      \token{Err(p)} & \token{v} is an option holding no value, whose error fits \token{p} & the variables of \token{p}\\[0pt]
      \token{T(f-> p, ...)} & \token{v} is an instance of \token{T} whose field \token{f} fits \token{p} & the variables of each \token{p}\\[0pt]
      \hline
    \end{tabular}
  \end{adjustbox}
\end{center}

An empty \token{Ok()} or \token{Err()} fits any option in that state
(Section~\ref{sec:compound:options}). The type \token{T} of a binding is optional,
and \token{\_} in its place stands for the type of the value; a type that is a
class tests the dynamic type of an instance, and binds \token{x} to it as a
\token{T}.

\begin{lstlisting}[style=coloredverbatim, caption=Patterns on integers]
fn describe(n: i32)-> [c8] {
  match n {
    -1 | 1 => "one, in absolute value"
    0 => "zero"
    2 ... 9 => "a digit"
    x if x < 0 => "negative"
    _ => "large"
  }
}

fn main() {
  println(describe(7));
}
\end{lstlisting}

\subsubsection{Variadic patterns}

A tuple or array pattern whose last element is followed by \token{...} fits a
value with at least as many elements as the pattern has before the dots. The
last pattern then fits the remaining elements, collected into a tuple for a
tuple, and into a slice for an array or a slice. With \token{\_} as the last
pattern, the remaining elements are ignored.

\begin{lstlisting}[style=coloredverbatim, caption=Variadic tuple and array patterns]
fn main() {
  let t = (1, 'c', 2.0);
  let (w, rest...) = t; // rest = ('c', 2.0), of type (c8, f64)
  println(w, " ", rest);

  let a = copy [1, 2, 3, 4];
  match a {
    [x, y, others...] => { // others = [3, 4], of type [i32]
      println(x, " ", y, " ", others);
    }
    _ => {
      println("fewer than two elements");
    }
  }
}
\end{lstlisting}

\subsubsection{Identifiers in patterns}

An identifier in a pattern is always a binding: it declares a new variable, and
never refers to an existing one. A binding whose name is already declared shadows
it, which is an error (E4199). To compare the value with a variable, the pattern
binds it under a new name and a guard compares the two.

\begin{lstlisting}[style=coloredverbatimError, escapechar=@]
fn check(speed: i32, limit: i32) {
  match speed {
    @\hb{limit}@ => { // error, declares a second /limit/
      println("at the limit");
    }
    s if s == limit => { // ok, compares with the parameter
      println("at the limit: ", s);
    }
    _ => {
      println("elsewhere");
    }
  }
}
\end{lstlisting}

A path such as \token{E::A} is not an identifier, and is compared with the
value like any expression; this is how an enumeration is matched.

\subsubsection{Decorators}

A binding can be decorated with \token{ref}, \token{mut} and \token{dmut}. The
value of the match must then give access to the matching memory, with
\token{ref} or \token{alias}, exactly as for a variable declared by a pattern
(Section~\ref{sec:memory:pattern_vdecl}), whose rules apply unchanged.

\subsection{Refutability}

A pattern is \textit{irrefutable} when it fits every value of its type, and
\textit{refutable} otherwise. The constructs that use patterns differ in which
they accept.

\begin{itemize}
\item A variable declaration requires an irrefutable pattern (E4193).
\item \token{if let} and \token{while let} require a refutable pattern (E4268,
  E4100).
\item \token{match} accepts both.
\end{itemize}

A pattern that can never fit, such as an array pattern of a length that the
matched array does not have, is an error in all four (E4006).

\subsection{Value and type}

A \token{match} is \textit{complete} when an arm is selected for every value:
either an arm fits every value -- \token{\_}, or a binding whose type, if
written, is the type of the value -- and has no guard, or a guard known at
compile time to hold; or the arms cover every member of an enumeration.
An arm that follows the point at which the match is complete can never be
selected, and is an error (E4259).

A complete match has a value, whose type is determined as for a conditional
(Section~\ref{sec:flow:cond_value_type}), with the bodies in the role of the
branches: breaking bodies are ignored, the type is that of the first remaining
body, and the others are implicitly cast to it. An incomplete match has no value,
and every body must be of type \token{void} (E4123). A \token{match} on a
\token{bool} whose arms are \token{true} and \token{false} is incomplete by this
rule, and needs \token{\_} to have a value.

\begin{lstlisting}[style=coloredverbatimError, escapechar=@]
fn name(day: i32)-> [c8] {
  match day {
    6 | 7 => @\hb{"weekend"}@ // error, the match has no default case
    1 ... 5 => @\hb{"weekday"}@
  }
}
\end{lstlisting}

\subsection{Diagnostics}

\begin{center}
  \begin{adjustbox}{max width=\linewidth}
    \begin{tabular}{|l|l|}
      \hline
      Code & Raised when\\[0pt]
      \hline
      \hline
      E4259 & an arm follows the point at which the match is complete\\[0pt]
      E4123 & an incomplete match has a body that produces a value\\[0pt]
      E4095 & two bodies have incompatible types\\[0pt]
      E4199 & a binding shadows an existing variable\\[0pt]
      E4006 & an array pattern has a length the matched array never has\\[0pt]
      E4193 & a variable declaration uses a refutable pattern\\[0pt]
      \hline
    \end{tabular}
  \end{adjustbox}
\end{center}
